\section{紧性的度量判据}
\label{sec:metric-compactness}

第~\ref{chap:topological-spaces}章把紧性定义为：每个开覆盖都有有限子覆盖。这个定义要求面对所有可能的开覆盖，直接核验往往很困难。在度量空间中可以换成两种更具体的检验：看点列是否总有收敛子列，或者看能否在任意精度下用有限个小球覆盖空间并保证极限存在。本节先证明前一种判据，再引入全有界性，把紧性拆成“有限尺度上可覆盖”与“逼近有极限”两个条件，最后给出子集相对紧的相应判据。

\subsection{紧性的序列判据}
\label{sec:metric-sequential-compactness}

闭区间上的有界点列总有收敛子列，这是卷一熟悉的事实。下面把这一性质作为一般度量空间的条件，并证明它恰好刻画紧性。

\begin{definition}[序列紧致]
\label{def:metric-sequentially-compact}
若度量空间$X$中的每个点列都有在$X$中收敛的子列，则称$X$是\term{序列紧致}（sequentially compact）的。
\end{definition}

证明紧性与序列紧致等价时，反方向需要一个把任意开覆盖与小球联系起来的引理。

\begin{lemma}[Lebesgue数]
\label{lem:metric-lebesgue-number}
设$X$序列紧致，$\{U_\alpha\}$是$X$的开覆盖。则存在$\delta>0$，使每个开球$B(x,\delta)$都含于某个$U_\alpha$。这样的$\delta$称为该覆盖的\term{Lebesgue数}（Lebesgue number）。
\end{lemma}

\begin{proof}
若不然，对每个$n$都有$x_n$使$B(x_n,1/n)$不含于任何$U_\alpha$。取收敛子列$x_{n_k}\to x$。$x$属于某个$U_\alpha$，取$r>0$使$B(x,r)\subseteq U_\alpha$。当$k$充分大时$d(x_{n_k},x)<r/2$且$1/n_k<r/2$，于是$B(x_{n_k},1/n_k)\subseteq B(x,r)\subseteq U_\alpha$，矛盾。
\end{proof}

\begin{theorem}[紧性的序列判据]
\label{thm:metric-compact-sequential}
度量空间$X$紧致，当且仅当它序列紧致。
\end{theorem}

\begin{proof}
先设$X$紧致，$(x_n)$是点列。假设它没有收敛子列。那么每个点$x$都有一个球$B(x,r_x)$，只含有限个下标$n$对应的$x_n$：否则对每个$k$，$B(x,1/k)$都含有无穷多项，可以依次选出$n_1<n_2<\cdots$使$x_{n_k}\in B(x,1/k)$，这个子列收敛到$x$。这些球覆盖$X$，取有限子覆盖，它们合起来只含有限个下标的项，但点列有无穷多个下标，矛盾。

再设$X$序列紧致，$\{U_\alpha\}$是开覆盖，取引理~\ref{lem:metric-lebesgue-number}的$\delta$。先说明$X$可以被有限个半径$\delta$的球覆盖。任取$y_1$；若$y_1,\ldots,y_k$已选好而$B(y_1,\delta),\ldots,B(y_k,\delta)$尚未覆盖$X$，就取一个未被覆盖的点作为$y_{k+1}$。若这个过程永不停止，得到的点列两两距离至少为$\delta$，它的任何子列都不是柯西列，因而不收敛，与序列紧致矛盾。所以有限步后得到覆盖$X=\bigcup_{j=1}^mB(y_j,\delta)$。每个$B(y_j,\delta)$含于某个$U_{\alpha_j}$，于是$U_{\alpha_1},\ldots,U_{\alpha_m}$是有限子覆盖。
\end{proof}

这个判据比较的是整个空间的性质，而不是某一个点列。$[0,1]$紧致，其中的点列$0,1,0,1,\ldots$并不收敛，但它有收敛子列，这并不妨碍紧性。反过来，$(0,1]$中的点列$1/n$在$\mathbb R$中收敛，但它的每个子列在$\mathbb R$中都趋于$0$，在$(0,1]$中没有极限，所以$(0,1]$不序列紧致，因而不紧。一个特定点列收敛与否，不能代替“每个点列都有收敛子列”。

在一般拓扑空间中，紧致与序列紧致互不蕴含，两个方向都有反例，这里不展开它们的构造。上面的证明在两处用到距离：正方向用半径$1/k$的可数个球选出子列，反方向用Lebesgue数和固定半径的球覆盖。

\subsection{全有界性}
\label{sec:metric-total-boundedness}

上面的证明中出现了一个独立于开覆盖的条件：任意给定半径，空间都能被有限个这样的球覆盖。这个条件描述在有限精度下空间是否只需要有限个代表。

\begin{definition}[全有界]
\label{def:metric-totally-bounded}
若对任意$\varepsilon>0$，存在有限个点$x_1,\ldots,x_m\in X$使$X=\bigcup_{j=1}^mB(x_j,\varepsilon)$，则称度量空间$X$是\term{全有界}（totally bounded）的。度量空间的子集$A$称为全有界的，是指子空间$(A,d_A)$全有界，即球心可以取在$A$中。
\end{definition}

第~\ref{sec:metric-finite-covers}节会说明，子集的球心改取在环境$X$中，只需把半径放大一倍，所以对子集而言两种约定得到同一个全有界概念。

全有界蕴含有界：取$\varepsilon=1$的有限覆盖，令$R=1+\max_jd(x_1,x_j)$，则$X\subseteq B(x_1,R)$。反过来不成立。设$X$是可数无限集，配上离散度量$d_0$。它有界，因为$\operatorname{diam}X=1$；它可分，因为$X$本身就是可数稠密集；但半径$1/2$的球都是单点集，覆盖无限集需要无穷多个，所以它不全有界。在$\mathbb R$中则不会出现这种差别：有界集含于某个区间$[-R,R]$，把区间等分成长度小于$\varepsilon$的有限段，每段至多需要一个以段内一点为中心的球，所以$\mathbb R$中的有界集都全有界。

全有界空间总是可分的。对每个$k$取有限个点组成的$1/k$覆盖中心集$D_k$，令$D=\bigcup_kD_k$，它可数。对任意$x$和$\varepsilon>0$，取$k>1/\varepsilon$，$x$落在某个$B(y,1/k)$中，$y\in D_k$，于是$d(x,y)<\varepsilon$，$D$稠密。上面离散空间的例子说明，可分不能反推全有界。

全有界还有一个序列刻画，它是连接覆盖与柯西列的关键。

\begin{theorem}[全有界的序列刻画]
\label{thm:metric-totally-bounded-cauchy}
度量空间$X$全有界，当且仅当$X$中每个点列都有柯西子列。
\end{theorem}

\begin{proof}
设$X$全有界，$(x_n)$是点列。用有限个半径$1$的球覆盖$X$，必有一个球含有点列的无穷多项，把这些项按原顺序记为子列$(x^{1}_n)$。再用有限个半径$1/2$的球覆盖$X$，其中一个球含有$(x^{1}_n)$的无穷多项，记为$(x^{2}_n)$。如此继续，第$k$个子列的全部项落在一个半径$1/k$的球中，并且每个子列都取自前一个。取对角子列$y_k=x^{k}_k$，它按原下标严格递增，是$(x_n)$的子列。当$k,l\geq K$时，$y_k,y_l$都属于第$K$个子列，落在同一个半径$1/K$的球中，$d(y_k,y_l)<2/K$，所以$(y_k)$是柯西列。

反过来，若对某个$\varepsilon>0$，$X$不能被有限个半径$\varepsilon$的球覆盖，按定理~\ref{thm:metric-compact-sequential}证明中的办法依次选点，得到两两距离至少为$\varepsilon$的点列，它没有柯西子列。
\end{proof}

\subsection{完备性与全有界性}
\label{sec:metric-complete-totally-bounded}

把序列判据拆开看，“每个点列都有收敛子列”包含两步：先选出彼此越来越近的子列，再保证它有极限。前一步由全有界负责，后一步由完备性负责。

\begin{theorem}[紧致的完备全有界判据]
\label{thm:metric-compact-complete-totally-bounded}
度量空间$X$紧致，当且仅当它完备且全有界。
\end{theorem}

\begin{proof}
设$X$紧致。全体半径$\varepsilon$的开球构成开覆盖，有限子覆盖说明$X$全有界。设$(x_n)$是柯西列，由定理~\ref{thm:metric-compact-sequential}它有收敛子列，再由引理~\ref{lem:metric-cauchy-basic}整列收敛，所以$X$完备。

设$X$完备且全有界。由定理~\ref{thm:metric-totally-bounded-cauchy}，每个点列都有柯西子列，由完备性该子列收敛，所以$X$序列紧致，从而紧致。
\end{proof}

两个条件缺一不可。$(0,1]$全有界但不完备，$1/n$这样的柯西列失去了极限；$\mathbb R$完备但不全有界，点列$x_n=n$两两距离至少为$1$，选不出柯西子列。

\begin{theorem}[欧氏空间中的紧集]
\label{thm:metric-heine-borel}
$(\mathbb R^n,\rho_\infty)$的子集$A$紧致，当且仅当$A$是有界闭集。
\end{theorem}

\begin{proof}
紧致度量空间全有界因而有界；它作为完备子空间，按定理~\ref{thm:metric-closed-complete}是闭集。反过来，设$A$有界且闭。$\mathbb R^n$完备，闭集$A$完备。有界性给出$A\subseteq[-R,R]^n$，把立方体按每个坐标分成长度小于$\varepsilon$的$k$段，得到$k^n$个小立方体，每个小立方体中两点的$\rho_\infty$距离小于$\varepsilon$。对每个与$A$相交的小立方体取$A$中的一点为球心，这有限个半径$\varepsilon$的球覆盖$A$，所以$A$全有界。由定理~\ref{thm:metric-compact-complete-totally-bounded}，$A$紧致。
\end{proof}

由于$\rho_\infty$、$\rho_1$与欧氏距离诱导相同拓扑，结论对欧氏距离同样成立。这个“有界闭集即紧集”的判据不能推广到任意度量空间。可数无限的离散空间作为自身是闭集，有界，也完备，却不全有界，所以不紧：它的开覆盖$\{\{x\}\}_{x\in X}$没有有限子覆盖。第~\ref{chap:normed-spaces}章还会看到，无限维赋范空间的闭单位球有界闭，同样不紧。在那里，起决定作用的仍是全有界性。

\subsection{相对紧性的判据}
\label{sec:metric-relative-compactness}

第~\ref{chap:topological-spaces}章把闭包紧致的子集称为相对紧集。在Hausdorff空间中，这等价于子集含于某个紧集：紧集是闭集，紧集的闭子集仍紧。下面给出度量空间中的序列刻画，并说明它与全有界的关系。

\begin{theorem}[相对紧的序列判据]
\label{thm:metric-relatively-compact-sequential}
设$A\subseteq X$。$A$相对紧，当且仅当$A$中每个点列都有在$X$中收敛的子列。
\end{theorem}

\begin{proof}
设$\overline A$紧致。$A$中的点列也是$\overline A$中的点列，由定理~\ref{thm:metric-compact-sequential}有子列收敛到$\overline A$中的点，自然在$X$中收敛。

反过来，设条件成立，$(y_n)$是$\overline A$中的点列。由定理~\ref{thm:metric-closure-sequential}，对每个$n$取$a_n\in A$使$d(a_n,y_n)<1/n$。按假设有子列$a_{n_k}\to x\in X$，则$d(y_{n_k},x)\leq d(y_{n_k},a_{n_k})+d(a_{n_k},x)\to0$。$x$是$A$中点列的极限，属于$\overline A$。所以$\overline A$序列紧致，从而紧致。
\end{proof}

判据只要求极限存在于环境$X$中，不要求属于$A$。在$\mathbb R$中，$A=(0,1)$相对紧，点列$1/(n+1)$的极限$0$就不在$A$中。环境同样要紧：若把$(0,1)$看作$X=(0,1)$自身的子集，它的闭包是$X$，不紧，所以它不是$X$的相对紧集。

\begin{theorem}[完备空间中的相对紧集]
\label{thm:metric-relatively-compact-totally-bounded}
设$X$完备，$A\subseteq X$。$A$相对紧，当且仅当$A$全有界。
\end{theorem}

\begin{proof}
设$\overline A$紧致，它全有界。对$\varepsilon>0$，用以$\overline A$中点为中心的有限个半径$\varepsilon/2$的球覆盖$\overline A$，舍去不与$A$相交的球，在其余每个球中取一点$a_j\in A$。若$a\in A$落在中心为$y_j$的球中，则$d(a,a_j)\leq d(a,y_j)+d(y_j,a_j)<\varepsilon$，所以以$a_j$为中心的半径$\varepsilon$的球覆盖$A$。

设$A$全有界。对$\varepsilon>0$，取$A$的有限个半径$\varepsilon/2$的球覆盖，中心为$a_1,\ldots,a_m$。$\overline A$中的点$y$与某个$a\in A$的距离小于$\varepsilon/2$，$a$又与某个$a_j$的距离小于$\varepsilon/2$，所以$\overline A\subseteq\bigcup_jB(a_j,\varepsilon)$，$\overline A$全有界。$\overline A$是完备空间的闭集，因而完备，由定理~\ref{thm:metric-compact-complete-totally-bounded}紧致。
\end{proof}

证明中前后两半分别说明：全有界的子集仍全有界，全有界集合的闭包仍全有界。环境的完备性只在最后一步使用，它不能省略。在$X=\mathbb Q$中，$A=\mathbb Q\cap[0,2]$全有界，但它在$\mathbb Q$中的闭包就是它自身，截断列$q_n$说明它不完备，所以不紧。没有环境完备性时，全有界只能保证：$\overline A$的完备化完备且全有界，因而紧致；这里$\overline A$的完备化是$[0,2]$。第~\ref{chap:function-class-covers}章判断函数类是否相对紧时，就是在完备的函数空间中核验全有界。
